% from ega1-1-fr.tex
\oldpage[I]{91}
Pour prouver que \textit{b}) entraîne \textit{d}\,1) et \textit{d}\,2),
supposons d'abord que $V=D(g)$ pour un $g\in A$, et que $\mathscr{F}$ soit
isomorphe à un faisceau $\widetilde{N}$ associé à un $A_g$-module $N$~; en
remplaçant $X$ par $V$ et $A$ par $A_g$
(\hyperref[I.1.3.6-fr]{1.3.6}), on peut se ramener au cas où $g=1$. Alors
$\Gamma(D(f),\widetilde{N})$ et $N_f$ s'identifient canoniquement
(\hyperref[I.1.3.6-fr]{1.3.6} et \hyperref[I.1.3.7-fr]{1.3.7}), donc une
section $s\in\Gamma(D(f),\widetilde{N})$ s'identifie à un élément de la
forme $z/f^n$, où $z\in N$~; la section $f^n s$ s'identifie à l'élément
$z/1$ de $N_f$, et est par suite restriction à $D(f)$ de la section de
$\widetilde{N}$ sur $X$ identifiée à l'élément $z\in N$~; d'où
\textit{d}\,1) dans ce cas. De même, $t\in\Gamma(X,\widetilde{N})$ est
identifié à un élément $z'\in N$, la restriction de $t$ à $D(f)$ est
identifiée à l'image $z'/1$ de $z'$ dans $N_f$, et dire que cette image est
nulle signifie qu'il existe $n\geq0$ tel que $f^n z'=0$ dans $N$, ou, ce
qui revient au même, $f^n t=0$.

Pour achever de prouver que \textit{b}) entraîne \textit{d}\,1) et
\textit{d}\,2), il suffira d'établir le lemme suivant~:

\begin{lemma}[1.4.1.1]
\label{I.1.4.1.1-fr}
Supposons que $V$ soit réunion finie d'ensembles de la forme $D(g_i)$, et
que chacun des faisceaux $\mathscr{F}|D(g_i)$,
$\mathscr{F}|(D(g_i)\cap D(g_j))=\mathscr{F}|D(g_i g_j)$ vérifie
\textit{d}\,1) et \textit{d}\,2)~; alors $\mathscr{F}$ possède les deux
propriétés suivantes~:
\begin{enumerate}
  \item[\textit{d}'\,1)] Pour tout $f\in A$ et toute section
  $s\in\Gamma(D(f)\cap V,\mathscr{F})$, il existe un entier $n\geq0$ tel
  que $f^n s$ se prolonge en une section de $\mathscr{F}$ sur $V$.
  \item[\textit{d}'\,2)] Pour tout $f\in A$ et toute section
  $t\in\Gamma(V,\mathscr{F})$ telle que la restriction de $t$ à
  $D(f)\cap V$ soit $0$, il existe un entier $n\geq0$ tel que $f^n t=0$.
\end{enumerate}
\end{lemma}

Prouvons d'abord \textit{d}'\,2)~: comme
$D(f)\cap D(g_i)=D(fg_i)$, il existe pour chaque $i$ un entier $n_i$ tel
que la restriction de $(fg_i)^{n_i}t$ à $D(g_i)$ soit nulle~: comme
l'image de $g_i$ dans $A_{g_i}$ est inversible, la restriction de
$f^{n_i}t$ à $D(g_i)$ est aussi nulle~; prenant pour $n$ le plus grand des
$n_i$, on a \textit{d}'\,2).

Pour démontrer \textit{d}'\,1), appliquons \textit{d}\,1) au faisceau
$\mathscr{F}|D(g_i)$~: il existe un entier $n_i\geq0$ et une section $s_i'$
de $\mathscr{F}$ sur $D(g_i)$ prolongeant la restriction de
$(fg_i)^{n_i}s$ à $D(fg_i)$~; comme l'image de $g_i$ dans $A_{g_i}$ est
inversible, il y a une section $s_i$ de $\mathscr{F}$ sur $D(g_i)$ telle
que $s_i'=g_i^{n_i}s_i$, et $s_i$ prolonge la restriction de $f^{n_i}s$ à
$D(fg_i)$~; on peut en outre supposer tous les $n_i$ égaux à un même
entier $n$. Par construction, la restriction de $s_i-s_j$ à
$D(f)\cap D(g_i)\cap D(g_j)=D(fg_i g_j)$ est nulle~; d'après
\textit{d}\,2) appliqué au faisceau $\mathscr{F}|D(g_i g_j)$, il existe un
entier $m_{ij}\geq0$ tel que la restriction à $D(g_i g_j)$ de
$(fg_i g_j)^{m_{ij}}(s_i-s_j)$ soit nulle~; comme l'image de $g_i g_j$
dans $A_{g_i g_j}$ est inversible, la restriction de
$f^{m_{ij}}(s_i-s_j)$ à $D(g_i g_j)$ est nulle. On peut alors supposer tous
les $m_{ij}$ égaux à un même entier $m$, et il existe donc une section
$s'\in\Gamma(V,\mathscr{F})$ prolongeant les $f^m s_i$~; cette section
prolonge par suite $f^{n+m}s$, d'où \textit{d}'\,1).

Reste à prouver que \textit{d}\,1) et \textit{d}\,2) entraînent
\textit{a}). Montrons d'abord que \textit{d}\,1) et \textit{d}\,2)
entraînent que ces conditions sont vérifiées pour tout faisceau
$\mathscr{F}|D(g)$, où $g\in A$ est tel que $D(g)\subset V$. C'est évident
pour \textit{d}\,1)~; d'autre part, si $t\in\Gamma(D(g),\mathscr{F})$ est
telle que sa restriction à $D(f)\subset D(g)$ soit nulle, il existe par
\textit{d}\,1) un entier $m\geq0$ tel que $g^m t$ se
